\section{Kernel PCA}

\mode<presentation>{
\begin{frame} 
    \begin{center} \huge
        \secname
    \end{center}
	\begin{center}
		\includegraphics[width=0.3\textwidth]{img/koffer}%
		
	\vspace{5mm}
	non-linear transformation $\Leftrightarrow$ kernel trick
	\end{center}
	
\end{frame}
}

\begin{frame}\frametitle{\secname}

We apply standard linear PCA on the \emph{transformed} version of the data
$
\big\{
\vec{\phi}_{(\vec{x}^{(\alpha)})}
\big\}_{\alpha=1}^{p}
$.

We will first assume we have $\vec{\phi}_{(\vec{x})}$ 
but we will eventually turn to $K_{\alpha \beta}$ 
which we can actually obtain.\\

\pause

\svspace{10mm}

\underline{Remark:}
\slidesonly{It might look different from the slides, but it's not.}
\notesonly{A difference between these notes and the lecture slides is that} 
\slidesonly{T}\notesonly{t}he lecture slides employ ``identity'' as the mapping. 
This is why you don't see $\phi$ in the derivations of Kernel PCA in the slides but rather see $\vec x$ used directly.\\

\end{frame}

\subsection{The Method}

\begin{frame}{Kernel PCA: Outline the method}

\mode<presentation>{
\begin{itemize}
\item center the ``immediate'' input to PCA

\pause
\only<2>{
\begin{center}
	\includegraphics[width=0.5\textwidth]{img/meme_realinput}%
\end{center}
}

\pause

\item compute covariance $\vec C_\phi$
\item solve the eigenvalue problem
\item normalize the eigenvectors
\item sort the eigenvectors and eigenvalues
\end{itemize}

Use Kernel PCA:

\begin{itemize}
\item projection
\item reconstruction

\end{itemize}

\pause

\only<4>{
\placeimage{9}{6}{img/meme_exactly}{width=5cm}%

\vspace{5mm}
The difference is in the details.
}

}

\end{frame}

\subsubsection{Centering the immediate input to PCA}

\begin{frame}{\subsubsecname}

Remember, we will first assume that we have the non-linear mapping $\phi$.\\

PCA assumes its input is centered.
Its direct input are the $\phi$'s. Therefore,
\begin{equation}
\frac{1}{p} \sum^{p}_{\alpha=1} \vec{\phi}_{(\vec{x}^{(\alpha)})} \eqexcl \vec 0
\end{equation}

\question{Isn't it enough to center $\vec X$?}

\pause

- No, Centering $\vec X$ does not guarantee it stays centered after the transformation.
Therefore, there is no need to center $\vec X$ beforehand.

Example:
\begin{equation}
\E[\vec x] = 0 \quad \nRightarrow \quad \E[\vec x^2] = 0
\end{equation}

\end{frame}

\subsubsection{The covariance matrix of the transformed data}

\begin{frame}{\subsubsecname}

Compute the covariance matrix $\vec C_{\phi}$ for $\vec{\phi}_{(\vec{x})}$:


\begin{equation} \label{eq:cov}
\vec C_{\phi} = \frac{1}{p} \sum_{\alpha=1}^{p} \vec{\phi}_{(\vec{x}^{(\alpha)})} \vec{\phi}^{\top}_{(\vec{x}^{(\alpha)})}
\end{equation}

\end{frame}

\subsubsection{The eigenvalue problem}

\begin{frame}{\subsubsecname}

Solve the eigenvalue problem:

\begin{equation} \label{eq:eig}
\vec C_{\phi} \, \vec e = \lambda \vec e
\end{equation}

Each eigenvector $\vec e_i$ with corresponding $\lambda_i \ne 0$ lies in the span of 
\begin{equation}
\big\{
\vec{\phi}_{(\vec{x}^{(\alpha)})}
\big\}_{\alpha=1}^{p}. \qquad \big(\vec{\phi}_{(\vec{x}^{(\alpha)})} = \vec{\phi}^{(\alpha)} \text{ for brevity}\big)
\end{equation}

\pause

Consequently, there exists a set of coefficients (i.e. a coefficient for each transformed observation)
\begin{equation}
\big\{
a^{(\alpha)}
\big\}_{\alpha=1}^{p}\,,
\end{equation} which satisfies the following:

\begin{equation}
\label{eq:ephi}
\vec e = \sum^{p}_{\beta=1} a^{(\beta)} 
%\vec{\phi}_{(\vec{x}^{(\beta)})}
\vec{\phi}^{(\beta)}
\end{equation}

\end{frame}

\mode<article>{

It is possible to arrive at the linear relationship in \eqref{eq:ephi} by substituting \eqref{eq:cov} into the eigenvalue problem in \eqref{eq:eig}:

\begin{align}
\underbrace{\frac{1}{p} \sum_{\alpha=1}^{p} \vec{\phi}^{(\alpha)} {\color{blue}\big(\vec{\phi}^{(\alpha)}\big)^\top} 
}_{=\,\vec C_{\phi}}
 \, 
{\color{blue}\vec e}
= \lambda \;\, \vec e\\
\intertext{with ${\color{blue}\big(\vec{\phi}^{(\alpha)}\big)^\top \vec e = \vec e^\top\vec{\phi}^{(\alpha)}}$ measuring the scalar projection $u_{(\vec{\phi}^{(\alpha)})}$ along $\vec e$:}
\frac{1}{p} \sum_{\alpha=1}^{p} u_{(\vec{\phi}^{(\alpha)})} \vec{\phi}^{(\alpha)}
= \lambda \;\, \vec e\\
\vec e = 
\frac{1}{p} \sum_{\alpha=1}^{p} \frac{u_{(\vec{\phi}^{(\alpha)})}}{\lambda} \vec{\phi}^{(\alpha)}\\
\intertext{This effectivley expresses $\vec e$ as a linear combination of the transformed observations:}
\vec e = 
\frac{1}{p} \sum_{\alpha=1}^{p} a^{(\alpha)} \vec{\phi}^{(\alpha)} \stackrel{\substack{\text{switch index}\\{\alpha \rightarrow \beta}}}{=} \frac{1}{p} \sum_{\beta=1}^{p} a^{(\beta)} \vec{\phi}^{(\beta)}
\end{align}

\paragraph{Deriving the transformed eigenvalue problem}

\eqref{eq:ephi} tells us that we can describe $\vec e$ in terms of the transformed observations (a weighted summation of $\phi$'s).
 The use of the index $\beta$ is only to avoid collisions with $\alpha$ later.
 
}

\begin{frame}{\subsubsecname}

\slidesonly{
\vspace{-12mm}
\hspace{7.8cm}
\StickyNote[1.7cm]{
	\begingroup
	\footnotesize
	\begin{equation} \label{eq:cov}
		\vec C_{\phi} = \frac{1}{p} \sum_{\alpha=1}^{p} \vec{\phi}^{(\alpha)} \big(\vec{\phi}^{(\alpha)}\big)^\top
	\end{equation}
	\begin{equation}
	\label{eq:ephi}
		\vec e = \sum^{p}_{\beta=1} a^{(\beta)} \vec{\phi}^{(\beta)}
	\end{equation}
	\endgroup
}[3.5cm]
\vspace{-22mm}
}

\notesonly{
Substituting \eqref{eq:cov} and \eqref{eq:ephi} into the eigenvalue problem \eqref{eq:eig}:
}
\slidesonly{
Express the eigenvalue problem in terms of $\phi$'s:

\begin{equation} \label{eq:eig}
\vec C_{\phi} \;\; \vec e \; = \; \lambda \;\; \vec e
\end{equation}

}

\begin{equation}
\underbrace{\frac{1}{p} \sum_{\alpha=1}^{p} \vec{\phi}^{(\alpha)} \big(\vec{\phi}^{(\alpha)}\big)^\top 
\vphantom{\sum^{p}_{\beta=1}}
}_{=\,\vec C_{\phi}}
 \, 
\underbrace{\sum^{p}_{\beta=1} a^{(\beta)} \vec{\phi}^{(\beta)}}_{=\,\vec e}
 = \lambda \;\,
\underbrace{\sum^{p}_{\beta=1} a^{(\beta)} \vec{\phi}^{(\beta)}}_{=\,\vec e}
\end{equation}

\pause

After rearranging the terms we get:
\begin{equation} \label{eq:eig2}
\frac{1}{p} \sum_{\alpha=1}^{p} \sum^{p}_{\beta=1} 
a^{(\beta)} \vec{\phi}^{(\alpha)}
\underbrace{
 \big(\vec{\phi}^{(\alpha)}\big)^\top  \,  \vec{\phi}^{(\beta)}
}_{\substack{\text{scalar product}\\ = K_{\alpha\beta}}}
 = \lambda 
\sum^{p}_{\beta=1} a^{(\beta)} \vec{\phi}^{(\beta)}
\end{equation}

We are one step closer to not needing the explicit mapping $\phi$.

\end{frame}

\begin{frame}{\subsubsecname}

\slidesonly{
\vspace{-12mm}
\hspace{7.8cm}
\StickyNote[1.7cm]{
	\begingroup
	\footnotesize
	\begin{equation} \label{eq:cov}
		\vec C_{\phi} = \frac{1}{p} \sum_{\alpha=1}^{p} \vec{\phi}^{(\alpha)} \big(\vec{\phi}^{(\alpha)}\big)^\top
	\end{equation}
	\begin{equation}
	\label{eq:ephi}
		\vec e = \sum^{p}_{\beta=1} a^{(\beta)} \vec{\phi}^{(\beta)}
	\end{equation}
	\endgroup
}[3.5cm]
\vspace{-22mm}
}

\notesonly{
Recall from \sectionref{sec:nonlin} that we are not even able to compute $\vec{\phi}_{(\vec{x})}$ but we now see it is possible to avoid the transformation altogether by exploiting the kernel trick (cf. \eqref{eq:trick}) by substituting 
$ K_{\alpha \beta} $ for
$
\vec{\phi}^{\top}_{(\vec{x}^{(\alpha)})}
 \, 
  \vec{\phi}_{(\vec{x}^{(\beta)})}
$

\eqref{eq:eig2} becomes:}

\begin{equation} \label{eq:eig3}
\frac{1}{p} \sum_{\alpha=1}^{p} \sum^{p}_{\beta=1} 
a^{(\beta)}
\vec{\phi}^{(\alpha)}
K_{\alpha \beta}
= \lambda 
\sum^{p}_{\beta=1} a^{(\beta)} \vec{\phi}^{(\beta)}
\slidesonly{\hspace{40mm}}
\end{equation}

We now proceed with reformulating the above until we no longer have any $\phi$'s:

\pause

\notesonly{We }left-multiply\notesonly{ \eqref{eq:eig3}} with $\big(\vec \phi_{(\vec x^{(\gamma)})}\big)^\top$, where $\gamma = 1, \ldots, p$.
 We can pull $\big(\vec \phi^{(\gamma)}\big)^\top$ directly into the sum on the \slidesonly{LHS}\notesonly{left-hand-side} and the sum on the \slidesonly{RHS}\notesonly{right-hand-side}:

\only<2,3>{
\begin{equation} \label{eq:eig4}
\frac{1}{p} \sum_{\alpha=1}^{p} \sum^{p}_{\beta=1} 
a^{(\beta)}
\underbrace{
\left(\vec \phi^{(\gamma)}\right)^\top
\vec{\phi}^{(\alpha)}
}_{=K_{\gamma \alpha}}
K_{\alpha \beta}
 = \lambda 
\sum^{p}_{\beta=1} a^{(\beta)} 
\underbrace{
\left(\vec \phi^{(\gamma)}\right)^\top \vec{\phi}^{(\beta)}
}_{=K_{\gamma \beta}}
\end{equation}
}

\pause

%\newpage

\notesonly{\eqref{eq:eig4} without the clutter:}

\begin{equation} \label{eq:eigK}
\frac{1}{p} \sum_{\alpha=1}^{p} \sum^{p}_{\beta=1} 
a^{(\beta)}
K_{\gamma \alpha}
K_{\alpha \beta}
 = \lambda 
\sum^{p}_{\beta=1} a^{(\beta)} 
K_{\gamma \beta} \quad {\only<4>{\color{red}}\forall \gamma}
\end{equation}

\slidesonly{
\only<4>{
Use matrix notation to reduce the clutter.
}
}



\end{frame}

\begin{frame}

\slidesonly{
\begin{equation}
\frac{1}{p} \sum_{\alpha=1}^{p} \sum^{p}_{\beta=1} 
a^{(\beta)}
K_{\gamma \alpha}
K_{\alpha \beta}
 = \lambda 
\sum^{p}_{\beta=1} a^{(\beta)} 
K_{\gamma \beta} \quad {\color{red}\forall \gamma}
\end{equation}
}

We want to solve \notesonly{\eqref{eq:eigK}}\slidesonly{the above} \textbf{for all} training samples $\gamma$.
We can further reduce the clutter by using matrix notation.\\ 
Specifically, by computing the \emph{kernel matrix} $\vec K=\{K_{\alpha\beta}\}$, where \\
\begin{equation}
K_{\alpha \beta} = 
k(\vec x^{(\alpha)}, \vec x^{(\beta)}) = 
\big(\vec{\phi}^{(\alpha)}\big)^\top 
		\vec{\phi}^{(\beta)}
\end{equation}

We end up with this formulation of the eigenvalue problem:

\begin{align}
	\frac{1}{p} \vec{K}^2 \vec{a} = \lambda \vec{K} \, \vec{a}\\
	\vec{K}^2 \vec{a} = p \lambda \vec{K} \, \vec{a}
\end{align}


\end{frame}

\begin{frame}{Side-note: Arriving at the matrix notation}

\notesonly{Side note: Going from \eqref{eq:eigK} to the matrix notation:\\}

\slidesonly{

\svspace{-3mm}

\begingroup
\footnotesize
\begin{equation} \label{eq:eigK}
\frac{1}{p} \sum_{\alpha=1}^{p} \sum^{p}_{\beta=1} 
a^{(\beta)}
K_{\gamma \alpha}
K_{\alpha \beta}
 = \lambda 
 \sum^{p}_{\beta=1} a^{(\beta)} 
K_{{\gamma} \beta} \quad {\forall \gamma}
\end{equation}
\endgroup

\svspace{-1mm}
}

First, we look at the RHS\notesonly{ of \eqref{eq:eigK}}:

\svspace{-5mm}

\begin{equation}
 \ldots = \lambda 
 \sum^{{\color{blue}p}}_{\color{blue}\beta=1} a^{({\color{blue}\beta})} 
K_{{\color{red}\gamma} {\color{blue}\beta}} \quad {\color{red}\forall \gamma}
\end{equation}

\begin{itemize}
\item ${\color{red}\gamma}$ corresponds to a specific row in the Kernel matrix $\vec K$.
\item Let $\vec a$ be a vector that stores the values $a^{({\color{blue}\beta})}$, ${\color{blue}\beta=1,\ldots,p}$.
\item ${\color{blue}\beta}$ iterates over the elements in the vector $\vec a$ and columns in $\vec K$.
\item If we switch terms to:
\svspace{-3mm}
\begin{equation}
 \ldots = \lambda 
 \sum^{{\color{blue}p}}_{\color{blue}\beta=1}  
K_{{\color{red}\gamma} {\color{blue}\beta}} \, a^{({\color{blue}\beta})} \quad {\color{red}\forall \gamma}
\label{eq:rearrangerhs}
\end{equation}

\notesonly{We see that the }RHS is essentially the inner product of the $\gamma$-th row with $\vec a$.
Recall that we're doing this ${\color{red}\forall \gamma}$ (i.e. all rows). Therefore RHS $\Rightarrow \lambda \vec K \, \vec a$ (a vector). 


\end{itemize}

\end{frame}

\definecolor{darkgreen}{rgb}{0,0.6,0}

\begin{frame}{Side-note: Arriving at the matrix notation - LHS}

Next, we look at the LHS\notesonly{ of \eqref{eq:eigK}}:

\svspace{-3mm}

\begin{equation}
 \frac{1}{p} \sum_{{\color{darkgreen}\alpha=1}}^{p} \sum^{p}_{{\color{blue}\beta=1}} 
a^{({\color{blue}\beta})}
K_{{\color{red}\gamma} {\color{darkgreen}\alpha}}
K_{{\color{darkgreen}\alpha} {\color{blue}\beta}} = \ldots
 \quad {\color{red}\forall \gamma}
\end{equation}

\svspace{-1mm}

\begin{itemize}
\item Same as before:
\begin{itemize}
\item ${\color{red}\gamma}$ corresponds to a specific row in the Kernel matrix $\vec K$.
\item Let $\vec a$ be a vector that stores the values $a^{({\color{blue}\beta})}$, ${\color{blue}\beta=1,\ldots,p}$.
\item ${\color{blue}\beta}$ iterates over the elements in the vector $\vec a$ and columns in $\vec K$.
\end{itemize}
\item ${\color{darkgreen}\alpha}$ iterates over the columns in $\vec K$.
\item $K_{{\color{red}\gamma} {\color{darkgreen}\alpha}}$ doesn't change with ${\color{blue}\beta}$, so we can move it out of the inner sum:
\svspace{-3mm}
\begin{align}
 \frac{1}{p} \sum_{{\color{darkgreen}\alpha=1}}^{p} 
 K_{{\color{red}\gamma} {\color{darkgreen}\alpha}}
 \underbrace{
\sum^{p}_{{\color{blue}\beta=1}} 
a^{({\color{blue}\beta})}
K_{{\color{darkgreen}\alpha} {\color{blue}\beta}}}_{\circledast} = \ldots
 \quad {\color{red}\forall \gamma}
\end{align}
\item $\circledast$ is the inner-product of the ${\color{darkgreen}\alpha}$-th row and $\vec a$.
\item LHS is essentially the inner product of the $\gamma$-th row with the \underline{vector} $\circledast$.
\item Recall ${\color{red}\forall \gamma}$. Therefore LHS $\Rightarrow \vec K \, \vec K \, \vec a = \vec K^2 \vec a$. \notesonly{$\vec K$ is a symmetric matrix} 

\end{itemize}

\end{frame}

\begin{frame}{Back to the eigenvalue problem}

\notesonly{Back to the eigenvalue problem:\\}

\slidesonly{
\begin{align}
	\vec{K}^2 \vec{a} = p \lambda \vec{K} \, \vec{a}
\end{align}
}

\pause

$\vec K$ appears on both sides. All the solutions that are of interest remain represented in 
the following simpler eigenvalue problem, which we refer to as the \emph{transformed eigenvalue problem}:
\begin{equation}
\label{eq:eigsimple1}
	\vec{K} \, \vec{a} = p \lambda \vec{a}
\end{equation}

\pause

We can interpret $\vec a$ as the \emph{eigenvector} of $\vec K$

\pause

\question{Do we really need $p$ in the transformed eigenvalue problem?}

\pause

No, by omitting the constant $p$ (optional), we can rely on finding solutions for $\lambda$ that absorb it:

\begin{equation}
\label{eq:eigsimple2}
	\vec{K} \, \vec{a} = \lambda \vec{a}
\end{equation}

\end{frame}

\begin{frame}

\question{What was all this about?}

All we've been doing so far is reformulate the eigenvalue problem such that we end up 
with a formulation that only contains terms of the inner product kernel.\\

\pause

\slidesonly{
\begin{minipage}{0.45\textwidth}
	\begin{center}
		\includegraphics[width=3cm]{img/meme_necessary}
    \end{center}
\end{minipage}
\pause
\begin{minipage}{0.45\textwidth}
	\begin{center}
		\includegraphics[width=4cm]{img/meme_ofcourse2}
    \end{center}
\end{minipage}
}

\notesonly{
\question{Was all this necessary?}

Yes, }because

\begin{enumerate}
\item we want to enable PCA to find non-linear correlations and
\item we don't have access to $\vec \phi_{(\vec x)}$.
\end{enumerate}

\end{frame}

\begin{frame}

\slidesonly{
\begin{center}
		\includegraphics[width=6cm]{img/meme_nowwhat}
\end{center}
}

\pause

Now that we've solved the eigenvalue problem, we continue with the remaining steps for PCA.

\pause


\slidesonly{
\begin{center}
		\includegraphics[width=2cm]{img/meme_right}
\end{center}
}

\end{frame}

\subsubsection{Normalize the eigenvectors}

\begin{frame}{\subsecname}

Before we can project anything onto the space spanned by the PCs, we need to know they are normalized first.

In linear PCA, we've relied on \textit{Eig} giving us normalized eigenvectors $\vec e$. 

\question{If we feed it with the kernel matrix will we get normalized $\vec a$?}

\pause

- Yes, Eig() gives us \underline{unit} vectors regardless.

\pause

\question{So, are we done talking about normalization?}

\pause

-No, because it's the $\vec e$'s that PCA needs normalized.

\end{frame}

\begin{frame}{\subsubsecname}

Recall, that PCA solves the eigenvalue problem

\begin{equation}
\vec C_{\phi} \, \vec e = \lambda \vec e
\end{equation}

At the solution, $\vec e^\top \vec e = 1$. $\vec e$ are normalized eigenvectors to use as PCs.

We don't have the $\phi$'s so we opted to solve the transformed eigenvalue problem instead:

\begin{equation}
%\label{eq:eigsimple1}
	\vec{K} \, \vec{a} = p \lambda \vec{a}
\end{equation}

At the solution, $\widetilde{\vec a}^\top \widetilde{\vec a} = 1 \quad \nRightarrow \quad \vec e^\top \vec e \ne 1$\\

\notesonly{Now that we are aware that we need further normalization, }we'll use $\widetilde{\vec a}_k$ where $k=1,\ldots,p$ to denote the \textbf{un}normalized eigenvectors from the transformed eigenvalue problem and $\vec a_k$ to denote that the vector has been normalized to use to correctly project onto the corresponding PC.

\end{frame}

\begin{frame}{\subsubsecname}

Recall\notesonly{ing \eqref{eq:ephi} (we add the index $k$ to denote which eigenvector):}
\begin{equation}
\label{eq:ephik}
\vec e_k = \sum^{p}_{\beta=1} a_k^{(\beta)} \vec{\phi}_{(\vec{x}^{(\beta)})},
\end{equation}

We want an expression for the norm $\vec e^{\top}_k \vec e_k$ that does not involve $\phi$'s. We left-multiply\slidesonly{ the above}\notesonly{ \eqref{eq:ephik}} with $\left(\vec e_k\right)^\top$:
\begin{align}
\vec e^{\top}_k \vec e_k &= \sum^{p}_{\alpha=1}  a_k^{(\alpha)} \vec{\phi}_{(\vec{x}^{(\alpha)})}^\top \sum^{p}_{\beta=1} a_k^{(\beta)} \vec{\phi}_{(\vec{x}^{(\beta)})} \\
&= \sum^{p}_{\alpha=1}  \sum^{p}_{\beta=1}  a_k^{(\beta)}  \underbrace{\vec{\phi}_{(\vec{x}^{(\alpha)})}^\top  \vec{\phi}_{(\vec{x}^{(\beta)})}} a_k^{(\alpha)} \\
&= \sum^{p}_{\alpha=1}  \sum^{p}_{\beta=1}  a_k^{(\beta)} \quad \; K_{\alpha\beta} \quad \; a_k^{(\alpha)} \\
&= \widetilde {\vec a}_k^\top \vec K \, \widetilde {\vec a}_k
\end{align}

\end{frame}

\begin{frame}{\subsubsecname}

\notesonly{
And when we plug \eqref{eq:eigsimple1} into the above:
}
\slidesonly{
From $\vec{K} \, \widetilde {\vec a}_k = p \lambda \widetilde {\vec a}_k$ follows:
}

\begin{equation}
\label{eq:eignorm}
\vec e^{\top}_k \vec e_k = 
\widetilde {\vec a}_k^\top 
\underbrace{p \lambda_k \, \widetilde {\vec a}_k}_{=\, \vec K \, \widetilde {\vec a}_k} 
= p \lambda_k \, \widetilde {\vec a}_k^\top \widetilde {\vec a}_k  \eqexcl 1
\end{equation}

\pause

w.l.o.g.

\svspace{-7mm}

\begin{equation}
\widetilde {\vec a}_k^\top \widetilde {\vec a}_k = 1
\end{equation}

Remark: a unit-length vector $\widetilde {\vec a}_k$ does not imply \emph{normlized} for PCA anymore.

\pause

\notesonly{
Scaling $\widetilde {\vec a}_k$ by $\frac{1}{\sqrt{p \lambda_k}}$ yields 
a vector in the same direction as $\widetilde {\vec a}_k$ to satisfy \notesonly{\eqref{eq:eignorm}}\slidesonly{$\vec e^{\top}_k \vec e_k \eqexcl 1$}.\\
}
With
\svspace{-5mm}
\begin{equation}
\vec a_k^{norm.} := \frac{1}{\sqrt{p \lambda_k}} \widetilde {\vec a}_k,
\end{equation}
follows:
\svspace{-5mm}
\begin{align}
\vec e^{\top}_k \vec e_k 
&= p \lambda_k \, \; \widetilde {\vec a}_k^\top \; \widetilde {\vec a}_k\\
&= p \lambda_k \, \left(\vec a_k^{\text{norm.}}\right)^\top \vec a_k^{\text{norm.}} \\
&= p \lambda_k \left(\frac{1}{\sqrt{p \lambda_k}} \widetilde {\vec a}_k\right)^\top \left(\frac{1}{\sqrt{p \lambda_k}} \widetilde {\vec a}_k\right)
= 1
\end{align}

\end{frame}

\subsubsection{Sorting}

\begin{frame}{\subsubsecname}

Sort the eigenvectors such that the corresponding eigenvalues are arranged in decreasing order. 

\slidesonly{
\begin{center}
		\includegraphics[width=4cm]{img/meme_sort}
\end{center}
}

\end{frame}

\subsubsection{Projection}

\begin{frame}{\subsubsecname}
\notesonly{In order t}\slidesonly{T}o project some observation $\vec x$ into the PC space, we first map it into \notesonly{the non-linear }space of $\phi$ 
and then project that into the space spanned by the PCs.\\

\pause

\slidesonly{
\begin{center}
		\includegraphics[width=4cm]{img/meme_mapx}
\end{center}
}

\notesonly{ 
By now we should expect that the transformation will be performed}\slidesonly{Transform} via the kernel trick.
\notesonly{We basically }represent this sample $\vec x$ by its relation to the \emph{training data} 
(i.e. the $p$ observations that were used to compute the PCs).

\end{frame}

\begin{frame}{\subsubsecname}

\slidesonly{
\visible<2->{
\vspace{-12mm}
\hspace{8.cm}
\StickyNote[1.cm]{
	\begingroup
	\footnotesize
	\begin{equation}
	\label{eq:ephi}
		\vec e = \sum^{p}_{\beta=1} a^{(\beta)} \vec{\phi}^{(\beta)}
	\end{equation}
	\endgroup
}[3.cm]
\vspace{-5mm}
}
}
 
\notesonly{We derive the projection for Kernel PCA by starting}\slidesonly{Start} with the projection used in linear PCA.

Recall in linear PCA, to get the component of $\vec x$ in the direction of the $k$-th PC we use:

\svspace{-5mm}

\begin{equation}
\label{eq:projlinx}
u_k(\vec x) = \vec e_k^\top \vec x
\end{equation}

Assuming we have a the $\phi$'s:

\svspace{-3mm}

\begin{equation}
\label{eq:projlin}
u_k(\phi_{(\vec x)}) = {\slidesonly{\color{magenta}}\vec e_k^\top} \vec \phi_{(\vec x)}
\end{equation}

\pause

\notesonly{
We substitute $\vec \phi_{(\vec x)}$ for $\vec x$ and plug \eqref{eq:ephi} into \eqref{eq:projlin}:
}

\visible<3>{

\svspace{-2mm}

\begin{align}
\label{eq:projk1}
u_k(\vec \phi_{(\vec x)}) &=
{\slidesonly{\color{magenta}}\sum^{p}_{\beta=1} a^{(\beta)}}
\hspace{-5mm}
\underbrace{
{\color{magenta}\vec{\phi}^\top_{(\vec{x}^{(\beta)})}} \, \vec \phi_{(\vec x)}
}_{\substack{
\text{recognize the familiar}\\
\text{scalar product?}
\\ = k(\vec x^{(\beta)}, \vec x) = K_{\beta,\vec x}}}\notesonly{\\
&} \kern-4mm = \sum_{\beta=1}^{p} a_k^{(\beta)} K_{\beta, \vec x}
\end{align}

\notesonly{Note that }$\vec x$ can be a sample that was used in computing the PCs or a completly new ``test'' point.
}
\end{frame}

\subsubsection{Reconstruction}

\begin{frame}{\subsubsecname}

\slidesonly{
\begin{center}
		\includegraphics[width=4cm]{img/meme_reconstruct}
\end{center}
}

Since we never had the transformation $\phi$ to begin with. 
It is not possible to simply project a sample from PC space back into the original $N$-dimensional input space. 
Algorithms exist that \emph{approximate} a ``pre-image'' of some new observation.

\end{frame}
